% Unit 073 English translation of chapter6.tex lines 321--556.
% Source: Methods of Algebra, Volume 2, commit
% 9a5803ff77dd3257484cb177f851a73770a59dd3 (CC BY 4.0).
% stable-unit-id: o014.aljabr2.chapter6.group-homology-cohomology
% Coverage: Complete Section 6.2, ``Group Homology and Cohomology'';
% stops before Section 6.3 at line 557.

% segment-id: o014.li2.chapter6.group-homology-cohomology.h001
\section{Group Homology and Cohomology}\label{sec:group-coh-sub}
\hypertarget{o014.aljabr2.chapter6.group-homology-cohomology}{ }
% segment-id: o014.li2.chapter6.group-homology-cohomology.p001
Throughout this section, the commutative ring $\Bbbk$ and the group $G$ remain
fixed.

% segment-id: o014.li2.chapter6.group-homology-cohomology.d001
\begin{definition}
	\index[sym1]{MGMG@$M^G$, $M_G$}
	For every $G$-module $M$, define the $\Bbbk$-modules
% segment-id: o014.li2.chapter6.group-homology-cohomology.q001
	\begin{align*}
		M^G & := \left\{ m \in M : \forall g \in G, \; gm = m \right\}, \\
		M_G & := M \big/ \lrangle{gm - m : g \in G, \; m \in M}.
	\end{align*}
	We call $M^G$ the submodule of \emph{invariants} of $M$, and $M_G$ the
	quotient module of \emph{coinvariants} of $M$.
\end{definition}

% segment-id: o014.li2.chapter6.group-homology-cohomology.p002
Taking invariants and coinvariants gives $\Bbbk$-linear functors
$G\dcate{Mod}\to\Bbbk\dcate{Mod}$. Both can be characterized as adjoints to
the inflation functor $\mathrm{Infl}$ of Definition~\ref{def:Res-Infl} in the
case $H=G$.

% segment-id: o014.li2.chapter6.group-homology-cohomology.pr001
\begin{proposition}\label{prop:inv-coinv-infl}
	Consider the functor
	$\mathrm{Infl}:=\mathrm{Infl}^G_{\{1\}}:\Bbbk\dcate{Mod}\to
	G\dcate{Mod}$, which sends each $\Bbbk$-module $N$ to the $G$-module with
	trivial action $gy=y$. There are two adjoint pairs
% segment-id: o014.li2.chapter6.group-homology-cohomology.g001
	\[\begin{tikzcd}[row sep=tiny]
		\mathrm{Infl}: \Bbbk\dcate{Mod} \arrow[shift left, r] & G\dcate{Mod} : (\cdot)^G , \arrow[shift left, l] \\
		(\cdot)_G: G\dcate{Mod} \arrow[shift left, r] & \Bbbk\dcate{Mod}: \mathrm{Infl}. \arrow[shift left, l]
	\end{tikzcd}\]
\end{proposition}
% segment-id: o014.li2.chapter6.group-homology-cohomology.pf001
\begin{proof}
	Let $M$ be a $G$-module and $N$ a $\Bbbk$-module. To specify a
	$G$-module homomorphism $\varphi:\mathrm{Infl}(N)\to M$ is the same as to
	specify a $\Bbbk$-module homomorphism $\varphi:N\to M$ such that
	$\varphi(y)=g\varphi(y)$ always holds, and this in turn is the same as to
	specify a $\Bbbk$-module homomorphism $\varphi:N\to M^G$.

	On the other hand, to specify a $\Bbbk$-module homomorphism
	$\psi:M_G\to N$ is the same as to specify a $\Bbbk$-module homomorphism
	$\varphi:M\to N$ such that $\varphi(gx-x)=0$ always holds. The latter
	condition is equivalent to $\varphi(gx)=g\varphi(x)$, where the $G$-action
	on the right is defined by the $G$-module structure $\mathrm{Infl}(N)$.
\end{proof}

% segment-id: o014.li2.chapter6.group-homology-cohomology.p003
The functors $(\cdot)^G$ and $(\cdot)_G$ can also be understood from the
viewpoint of $\Bbbk[G]$-modules. In what follows, through the augmentation
homomorphism of Remark~\ref{rem:augmentation-kG}, we regard $\Bbbk$ as a
$(\Bbbk[G],\Bbbk[G])$-bimodule; in other words, the action of $G$ on either
side of $\Bbbk$ is trivial.

% segment-id: o014.li2.chapter6.group-homology-cohomology.pr002
\begin{proposition}\label{prop:inv-coinv-Hom}
	Upon identifying $G$-modules with left $\Bbbk[G]$-modules, there are natural
	isomorphisms of functors
% segment-id: o014.li2.chapter6.group-homology-cohomology.q002
	\begin{gather*}
		(\cdot)^G \simeq \Hom_G(\Bbbk, \cdot), \quad  (\cdot)_G \simeq \Bbbk \dotimes{\Bbbk[G]} (\cdot).
	\end{gather*}
\end{proposition}
% segment-id: o014.li2.chapter6.group-homology-cohomology.pf002
\begin{proof}
	Consider a $G$-module $M$. By Proposition~\ref{prop:inv-coinv-infl}, to
	specify a $G$-module homomorphism $\varphi:\Bbbk\to M$ is the same as to
	specify a $\Bbbk$-module homomorphism $\varphi:\Bbbk\to M^G$. The latter
	is equivalent to specifying the element $\varphi(1)$ of $M^G$. This gives
	the first isomorphism.

	For the second isomorphism, first identify $\Bbbk$ with
	$\Bbbk[G]/\mathfrak{I}$, where the augmentation ideal is
	$\mathfrak{I}=\bigoplus_{g\neq1_G}\Bbbk(g-1_G)$. Then
	$\Bbbk\dotimes{\Bbbk[G]}M\simeq M/\mathfrak{I}M=M_G$, by the map sending
	$t\otimes x$ to the image of $tx\in M$ in $M_G$. This proves the claim.
\end{proof}

% segment-id: o014.li2.chapter6.group-homology-cohomology.p004
Thus the adjunctions of Proposition~\ref{prop:inv-coinv-infl} can be understood
as the special case of Proposition~\ref{prop:pullback-two-adjoint} obtained by
taking the trivial homomorphism $\varphi:G\twoheadrightarrow\{1\}$.

% segment-id: o014.li2.chapter6.group-homology-cohomology.co001
\begin{corollary}\label{prop:Ind-invariant}
	For every group homomorphism $\varphi:H\to G$, there are isomorphisms of
	functors
% segment-id: o014.li2.chapter6.group-homology-cohomology.q003
	\[\begin{array}{rlrl}
		\Hom_H(\Bbbk[G], \cdot)^G & \rightiso (\cdot)^H, & (\cdot)_H & \rightiso (\Bbbk[G] \dotimes{\Bbbk[H]} \cdot)_G , \\
		\varphi & \mapsto \varphi(1_G), & (x \;\text{in the quotient}) & \mapsto (1_G \otimes x \;\text{in the quotient}) .
	\end{array}\]
	Both sides are functors from $H\dcate{Mod}$ to $\Bbbk\dcate{Mod}$. The
	left-hand side involves the functor
	$\Hom_H(\Bbbk[G],\cdot):H\dcate{Mod}\to G\dcate{Mod}$ introduced in
	Proposition~\ref{prop:pullback-two-adjoint}.
\end{corollary}
% segment-id: o014.li2.chapter6.group-homology-cohomology.pf003
\begin{proof}
	Let $N$ be an $H$-module. The first isomorphism follows from the adjunction
	in Proposition~\ref{prop:pullback-two-adjoint}:
% segment-id: o014.li2.chapter6.group-homology-cohomology.q004
	\[ \Hom_G(\Bbbk, \Hom_H(\Bbbk[G], N)) \simeq \Hom_H(\varphi^* \Bbbk, N) = \Hom_H(\Bbbk, N). \]

	The second isomorphism follows from the associativity constraint for tensor
	products,
	$\Bbbk\dotimes{\Bbbk[G]}(\Bbbk[G]\dotimes{\Bbbk[H]}N)
	\simeq\Bbbk\dotimes{\Bbbk[H]}N$. Verification of the explicit map is left
	to the reader. Alternatively, the assertion can be proved directly by
	showing that the displayed map is well defined and invertible.
\end{proof}

% segment-id: o014.li2.chapter6.group-homology-cohomology.p005
The adjoint pairs in Proposition~\ref{prop:inv-coinv-infl} also ensure that
$(\cdot)_G$ is right exact and $(\cdot)^G$ is left exact; these facts are also
easy to verify directly.

% segment-id: o014.li2.chapter6.group-homology-cohomology.d002
\begin{definition}[Group homology and cohomology]\label{def:group-coh-ho}
	\index{group homology and group cohomology}
	\index[sym1]{HnGM@$\Hm_n(G, M)$, $\Hm^n(G, M)$}
	Let $G$ be a group. \emph{Group homology} with coefficients in $\Bbbk$ is
	defined to be the left-derived functors of $(\cdot)_G$:
% segment-id: o014.li2.chapter6.group-homology-cohomology.q005
	\[ \Hm_n(G, \cdot) := \mathrm{L}_n (\cdot)_G: G\dcate{Mod} \to \Bbbk\dcate{Mod}, \]
	while \emph{group cohomology} is defined to be the right-derived functors of
	$(\cdot)^G$:
% segment-id: o014.li2.chapter6.group-homology-cohomology.q006
	\[ \Hm^n(G, \cdot) := \mathrm{R}^n (\cdot)^G: G\dcate{Mod} \to \Bbbk\dcate{Mod}, \]
	where $n\in\Z$. Both vanish unless $n\geq0$. Similar definitions apply to
	right $G$-modules.
\end{definition}

% segment-id: o014.li2.chapter6.group-homology-cohomology.p006
Following the convention of much of the literature, when the coefficients are
not specified we assume $\Bbbk=\Z$. In this case, a $G$-module is simply an
additive group equipped with a left $G$-action that preserves addition.

% segment-id: o014.li2.chapter6.group-homology-cohomology.e001
\begin{example}\label{eg:group-H1}
	Consider the augmentation ideal
	$\mathfrak{I}=\bigoplus_{g\neq1_G}\Bbbk(g-1_G)$ of $\Bbbk[G]$. For every
	$G$-module $M$, there is an isomorphism
	$\Hm_0(G,M)\simeq M_G\simeq M/\mathfrak{I}M$; in particular,
	$\Hm_0(G,\mathfrak{I})\simeq\mathfrak{I}/\mathfrak{I}^2$.

	Now consider the short exact sequence
	$0\to\mathfrak{I}\to\Bbbk[G]\to\Bbbk\to0$, whose final morphism is the
	augmentation homomorphism. Since $\Bbbk[G]$ is a projective module, the
	dimension-shifting technique of Proposition~\ref{prop:dimension-shifting}
	gives
% segment-id: o014.li2.chapter6.group-homology-cohomology.q007
	\[ \Hm_n(G, \Bbbk) \simeq \begin{cases}
		\Hm_{n-1}(G, \mathfrak{I}), & n \geq 2 \\
		\Ker\left[ \mathfrak{I}/\mathfrak{I}^2 \to \Bbbk[G]/\mathfrak{I} \right] = \mathfrak{I}/\mathfrak{I}^2 , & n = 1.
	\end{cases}\]
\end{example}

% segment-id: o014.li2.chapter6.group-homology-cohomology.p007
By the isomorphism between $\Bbbk[G]\dcate{Mod}$ and $G\dcate{Mod}$, together
with Proposition~\ref{prop:inv-coinv-Hom}, group homology and cohomology can
also be interpreted through the functors $\Tor$ and $\Ext$ introduced in
\S\ref{sec:Ext-Tor}:
% segment-id: o014.li2.chapter6.group-homology-cohomology.q008
\begin{equation}\label{eqn:grp-Hm-Tor-Ext}\begin{aligned}
	\Hm_n(G, M) & \simeq \Tor^G_n(\Bbbk, M) := \Tor^{\Bbbk[G]}_n(\Bbbk, M), \\
	\Hm^n(G, M) & \simeq \Ext_G^n(\Bbbk, M) := \Ext_{\Bbbk[G]}^n(\Bbbk, M).
\end{aligned}\end{equation}

% segment-id: o014.li2.chapter6.group-homology-cohomology.p008
This interpretation sometimes reduces the questions at hand to known results
in module theory. Here is a fundamental, nontrivial example involving a
spectral sequence relating $\Ext$ and $\Tor$.

% segment-id: o014.li2.chapter6.group-homology-cohomology.pr003
\begin{proposition}[Universal coefficient theorem for group homology and cohomology]\label{prop:group-UCT}
	\index{universal coefficient theorem}
	Take $\Bbbk=\Z$. Equip the $\Z$-module $M$ with the trivial $G$-action,
	making it a $G$-module. For every $n\in\Z_{\geq0}$, there are canonical
	short exact sequences
% segment-id: o014.li2.chapter6.group-homology-cohomology.g002
	\[\begin{tikzcd}[row sep=tiny, column sep=small]
		0 \arrow[r] & \Ext^1_{\Z}\left( \Hm_{n-1}(G, \Z), M \right) \arrow[r] & \Hm^n(G, M) \arrow[r] & \Hom_{\Z}\left( \Hm_n(G, \Z), M \right) \arrow[r] & 0, \\
		0 \arrow[r] & M \dotimes{\Z} \Hm_n(G, \Z) \arrow[r] & \Hm_n(G, M) \arrow[r] & \Tor^{\Z}_1(M, \Hm_{n-1}(G, \Z)) \arrow[r] & 0;
	\end{tikzcd}\]
	for $n=0$, use the convention $\Hm_{-1}(G,\Z)=0$.
\end{proposition}
% segment-id: o014.li2.chapter6.group-homology-cohomology.pf004
\begin{proof}
	We may assume $n\geq1$. First consider cohomology. Once the reader is
	familiar with the standard cochain and chain complexes to be introduced
	below, the short exact sequence follows easily by applying the universal
	coefficient theorem for cohomology from the exercises in \CHref{sec:cplx}
	to the chain complex $C:=\mathsf{L}\dotimes{\Bbbk[G]}\Z$ and the
	$\Z$-module $M$. The details are left to interested readers.

	We give another argument here. In Example~\ref{eg:change-of-rings-SS-bis},
	take the ring homomorphism $R:=\Z[G]\to S:=\Z$ to be the augmentation;
	also take the left $R$-module $Y=\Z$ (the trivial $G$-module) and the left
	$S$-module $X=M$. The short exact sequence in that example then becomes
	the required one.

	For homology, one may similarly apply the standard chain complex to the
	homological K\"unneth theorem~\ref{prop:Kunneth-homology}, leaving the
	details to the reader, or understand the result through
	Example~\ref{eg:change-of-rings-SS-bis}.
\end{proof}

% segment-id: o014.li2.chapter6.group-homology-cohomology.p009
Following the constructions introduced in \S\ref{sec:derived-primer}, we may
also take the relevant derived functors of a bounded-below chain complex (or a
bounded-below cochain complex) $M$. These are called group hyperhomology (or
group hypercohomology). In the derived category one may further consider
% segment-id: o014.li2.chapter6.group-homology-cohomology.q009
\[ \Bbbk \otimesL[{\Bbbk[G]}] M, \quad \RHom_G(\Bbbk, M) := \RHom_{\Bbbk[G]}(\Bbbk, M). \]

% segment-id: o014.li2.chapter6.group-homology-cohomology.p010
By \eqref{eqn:grp-Hm-Tor-Ext} and the fact that $\Ext$ (respectively $\Tor$)
can be computed using a projective resolution in the first variable, the free
resolution $\mathsf{L}\to\Bbbk$ of
Proposition~\ref{prop:trivial-mod-resolution} immediately supplies a concrete
complex (respectively chain complex) for computing group cohomology
(respectively homology). Definition~\ref{def:resolution-trivial-module} gives
two ways to make $\mathsf{L}$ a complex of left or right $G$-modules: use
$\partial_n:\mathsf{L}_n\to\mathsf{L}_{n-1}$ or use
$\partial'_n:\mathsf{L}_n\to\mathsf{L}_{n-1}$. We shall choose between them
according to the context so as to obtain the standard form of the complex;
the choice does not affect its cohomology.

% segment-id: o014.li2.chapter6.group-homology-cohomology.pr004
\begin{proposition}[Computing group cohomology by the standard cochain complex]
\label{prop:groupcoh-cochain}
	\index{standard cochain complex}
	\index[sym1]{CGM@$C(G, M)$}
	For every $G$-module $M$, define the \emph{standard cochain complex}
	$C(G,M)=(C^n(G,M))_n$ of $\Bbbk$-modules as follows:
% segment-id: o014.li2.chapter6.group-homology-cohomology.q010
	\begin{gather*}
		C^n(G, M) := \left\{ \text{maps}\; f: G^n \to M \right\}, \quad n \in \Z_{\geq 0},
	\end{gather*}
	Each term is made into a $\Bbbk$-module by pointwise addition and scalar
	multiplication, the terms in negative degrees are defined to be $0$, and
	$d^n:C^n(G,M)\to C^{n+1}(G,M)$ is defined by
% segment-id: o014.li2.chapter6.group-homology-cohomology.q011
	\begin{multline*}
		(d^n f)(g_1, \ldots, g_{n+1}) = g_1 f(g_2, \ldots, g_{n+1}) \\
		+ \sum_{k=1}^n (-1)^k f(\ldots, g_k g_{k+1}, \ldots) + (-1)^{n+1} f(g_1, \ldots, g_n).
	\end{multline*}

	For every $n$, there is a canonical isomorphism
	$\Hm^n(C(G,M))\simeq\Hm^n(G,M)$. More precisely, $C(G,M)$ represents
	$\RHom_G(\Bbbk,M)$ in the derived category.
\end{proposition}
% segment-id: o014.li2.chapter6.group-homology-cohomology.pf005
\begin{proof}
	Consider the complex $\mathsf{L}=(\mathsf{L}_n,\partial_n)_n$. Since
	$\mathsf{L}\to\Bbbk$ is a free resolution by left $G$-modules,
	$\RHom_G(\Bbbk,M)$ is represented by the following $\Hom$ complex
	(Definition~\ref{def:Hom-cplx}):
% segment-id: o014.li2.chapter6.group-homology-cohomology.q012
	\begin{equation*}
		\Hom^\bullet(\mathsf{L}, M) = \left[ \cdots \to \Hom_{\Bbbk[G]}(\mathsf{L}_n, M) \xrightarrow{(-1)^{n+1} \partial_{n+1}^*} \Hom_{\Bbbk[G]}(\mathsf{L}_{n+1}, M) \to \cdots \right] ,
	\end{equation*}
	Its degrees are $\ldots,-n,-n-1,\ldots$; the sign $(-1)^{n+1}$ comes
	from the general definition of the $\Hom$ complex.

	For every $n\geq0$, use the basis $(g_1|\cdots|g_n)$ of the free left
	$\Bbbk[G]$-module $\mathsf{L}_n$ from \eqref{eqn:L-basis}. Thus
% segment-id: o014.li2.chapter6.group-homology-cohomology.g003
	\[\begin{tikzcd}[row sep=tiny, column sep=tiny]
		\Hom_G(\mathsf{L}_n, M) \arrow[r, "\sim"] & C^n(G, M) \\
		\phi \arrow[mapsto, r] & {\left[ f: (g_1, \ldots, g_n) \mapsto \phi\left( (g_1 | \cdots | g_n) \right)\right]} ,
	\end{tikzcd}\]
	which is plainly an isomorphism of $\Bbbk$-modules. The key point is to
	identify $(-1)^{n+1}\partial_{n+1}^*$ on $C^n(G,M)$. Consider
	$(-1)^{n+1}\partial_{n+1}^*\phi=(-1)^{n+1}\phi\circ\partial_{n+1}$.
	By \eqref{eqn:L-basis-partial}, this map sends
	$(g_1|\cdots|g_{n+1})\in\mathsf{L}_{n+1}$ to
% segment-id: o014.li2.chapter6.group-homology-cohomology.q013
	\begin{multline*}
		\phi\left( g_1 (g_2 | \cdots | g_{n+1})\right) + \sum_{k=1}^n (-1)^k \phi\left((\cdots | g_k g_{k+1}| \cdots)\right) + \cdots + (-1)^{n+1} \phi\left((g_1 | \cdots |g_n)\right) \\
		= g_1 f(g_2, \ldots, g_{n+1}) + \sum_{k=1}^n (-1)^k f(\ldots, g_k g_{k+1}, \ldots) + (-1)^{n+1} f(g_1, \ldots, g_n).
	\end{multline*}
	This is precisely $(d^nf)(g_1,\ldots,g_{n+1})$ in the
	statement.\footnote{Translator's correction: the source ends the argument
	list at $g_n$, but $d^nf\in C^{n+1}(G,M)$ and the calculation above uses
	$(g_1,\ldots,g_{n+1})$.}
\end{proof}

% segment-id: o014.li2.chapter6.group-homology-cohomology.p011
The cohomology of a right $G$-module $M$ has a similar description, which we
do not repeat. We next state the corresponding result for the group homology
of a left $G$-module; its proof is immediate.

% segment-id: o014.li2.chapter6.group-homology-cohomology.pr005
\begin{proposition}[Computing group homology by the standard chain complex]
\label{prop:grouph-chain}
	For every $G$-module $M$ and $n\in\Z_{\geq0}$, there is a canonical
	isomorphism
% segment-id: o014.li2.chapter6.group-homology-cohomology.q014
	\[ \Hm_n(G, M) \simeq \Hm_n\left( \mathsf{L} \dotimes{\Bbbk[G]} M\right), \quad n \in \Z_{\geq 0}, \]
	where $(\mathsf{L}_n,\partial'_n)_n$ is used as the chain complex of right
	$G$-modules. Take the basis $(g_1|\cdots|g_n)^\dagger$ of the free right
	$\Bbbk[G]$-module $\mathsf{L}_n$ from \eqref{eqn:L-basis-right}. Then
	$\mathsf{L}_n\dotimes{\Bbbk[G]}M$ is identified with the direct
	sum\footnote{Translator's correction: the source prints $\mathsf{L}^n$,
	whereas the chain complex and the basis in the same sentence use
	$\mathsf{L}_n$.} $M^{\oplus G^n}$, and an element in the summand
	corresponding to $(g_1,\ldots,g_n)\in G^n$ has a unique expression
% segment-id: o014.li2.chapter6.group-homology-cohomology.q015
	\[ (g_1 | \cdots | g_n)^\dagger \otimes x, \quad x \in M, \]
	while, by \eqref{eqn:L-basis-partial-right},
	$\partial'_n\otimes\identity_M:\mathsf{L}_n\dotimes{\Bbbk[G]}M\to
	\mathsf{L}_{n-1}\dotimes{\Bbbk[G]}M$ is given by
% segment-id: o014.li2.chapter6.group-homology-cohomology.q016
	\begin{multline*}
		(g_1 | \cdots | g_n)^\dagger \otimes x \mapsto \\
		(g_2 | \cdots | g_n)^\dagger \otimes x + \sum_{k=1}^{n-1} (-1)^k (\cdots | g_k g_{k+1} | \cdots)^\dagger \otimes x \\
		+ (-1)^n (g_1 | \cdots | g_{n-1})^\dagger \otimes g_n x.
	\end{multline*}

	For comparison with Proposition~\ref{prop:groupcoh-cochain}, denote the
	chain complex above by $(C_n(G,M))_n$ and call it the \emph{standard chain
	complex}. In the derived category it represents
	$\Bbbk\otimesL[{\Bbbk[G]}]M$.
\end{proposition}

% segment-id: o014.li2.chapter6.group-homology-cohomology.p012
Example~\ref{eg:group-H-comonad} below will explain the relation between the
standard chain complex and the bar construction.

% segment-id: o014.li2.chapter6.group-homology-cohomology.p013
There is of course a corresponding description for the homology of a right
$G$-module $M$. Consider $M\dotimes{\Bbbk[G]}\mathsf{L}$. We have
$M\dotimes{\Bbbk[G]}\mathsf{L}_n\simeq M^{\oplus G^n}$, and an element in
the summand corresponding to $(g_1,\ldots,g_n)$ has a unique expression
% segment-id: o014.li2.chapter6.group-homology-cohomology.q017
\[ x \otimes (g_1 | \cdots | g_n), \quad x \in M, \]
while the map $\identity_M\otimes\partial'_n$ is
% segment-id: o014.li2.chapter6.group-homology-cohomology.q018
\begin{multline*}
	x \otimes (g_1 | \cdots | g_n) \mapsto \\
	x g_1 (g_2 | \cdots | g_n) + \sum_{k=1}^{n-1} (-1)^k x \otimes (\cdots | g_k g_{k+1} | \cdots) \\
	+ (-1)^n x \otimes (g_1 | \cdots | g_{n-1}).
\end{multline*}
The corresponding standard chain complex is again denoted by
$(C_n(G,M))_n$.

% segment-id: o014.li2.chapter6.group-homology-cohomology.p014
Readers familiar with \S\ref{sec:HH} will readily recognize these complexes
as special cases of the Hochschild cochain and chain complexes. Here is the
precise statement.

% segment-id: o014.li2.chapter6.group-homology-cohomology.pr006
\begin{proposition}[Relation with Hochschild theory]
	Let $R:=\Bbbk[G]$. For a left (respectively right) $G$-module $M$, write
	$M_+$ (respectively ${}_+M$) for the $(R,R)$-bimodule determined by
% segment-id: o014.li2.chapter6.group-homology-cohomology.q019
	\[ g_1 x g_2 := g_1 x \; (\text{respectively}\; g_1 x g_2 = x g_2 ), \quad g_1 , g_2 \in G, \; x \in M. \]
	There are canonical isomorphisms compatible with long exact sequences,
% segment-id: o014.li2.chapter6.group-homology-cohomology.q020
	\[ \Hm^n(G, M) \simeq \HHm^n(R, M_+), \quad (\text{respectively}\; \Hm_n(G, M) \simeq \HHm_n(R, {}_+ M)), \]
	and in fact the complex $(C^n(G,M))_n$ (respectively the chain complex
	$(C_n(G,M))_n$) is isomorphic to $(C^n(R,M_+))_n$ (respectively
	$(C_n(R,{}_+M))_n$) from \eqref{eqn:HH-cplx}.
\end{proposition}
% segment-id: o014.li2.chapter6.group-homology-cohomology.pf006
\begin{proof}
	Observe that $R^{\otimes n}$ is a free $\Bbbk$-module with basis $G^n$:
	$(g_1,\ldots,g_n)$ corresponds to
	$g_1\otimes\cdots\otimes g_n\in R^{\otimes n}$. It remains only to
	compare the standard cochain complex with the complex described in
	\eqref{eqn:HH-cplx}. Concretely, in the cohomological case, send
	$f:G^n\to M$ to the corresponding $n$-multilinear $\Bbbk$-map $R^n\to M$.
	In the homological case, send
	$x\otimes(g_1|\cdots|g_n)^\dagger$ to
	$(x|g_1|\cdots|g_n)\in C_n(R,{}_+M)$. The remaining comparison is left to
	the reader.
\end{proof}

% segment-id: o014.li2.chapter6.group-homology-cohomology.p015
Returning to our main subject, from now on all $G$-modules under consideration
are left modules unless otherwise stated.

% segment-id: o014.li2.chapter6.group-homology-cohomology.e002
\begin{example}\label{eg:group-H1-further}
	In light of the discussion above, consider
% segment-id: o014.li2.chapter6.group-homology-cohomology.q021
	\[ \Hm_1(G, \Bbbk) \simeq \frac{\Ker\left[ C_1(G, \Bbbk) \to C_0(G, \Bbbk) \right]}{\Image\left[ C_2(G, \Bbbk) \to C_1(G, \Bbbk) \right]}. \]

	Every element of $C_1(G,\Bbbk)$ has a unique expression as a
	$\Bbbk$-linear combination of the elements
	$[g]:=(g)^\dagger\otimes1$ for $g\in G$. Observe that
	$C_0(G,\Bbbk)\simeq\Bbbk$ and
% segment-id: o014.li2.chapter6.group-homology-cohomology.q022
	\[ \left( \partial'_1 \otimes \identity_{\Bbbk}\right)[g] = 1 - 1 = 0, \quad \left( \partial'_2 \otimes \identity_{\Bbbk}\right)((g_1 | g_2)^\dagger \otimes 1) = [g_2] - [g_1 g_2] + [g_1]. \]
	Therefore $\Hm_1(G,\Bbbk)$ is isomorphic to the quotient $\Bbbk$-module
% segment-id: o014.li2.chapter6.group-homology-cohomology.q023
	\[ Q := \bigoplus_{g \in G} \Bbbk \cdot [g] \bigg/ \lrangle{ [g_1 g_2] - [g_1] + [g_2] : g_1, g_2 \in G }; \]
	The map $g\mapsto[g]$ induces a group homomorphism $a:G\to(Q,+)$. It has
	the following universal property: for every $\Bbbk$-module $N$ and group
	homomorphism $b:G\to(N,+)$, there is a unique $\Bbbk$-module homomorphism
	$\varphi:Q\to N$ such that $b=\varphi a$. The reader may quickly verify
	this universal property and deduce from it the canonical isomorphism of
	$\Bbbk$-modules
% segment-id: o014.li2.chapter6.group-homology-cohomology.g004
	\[\begin{tikzcd}[row sep=tiny]
		G_{\mathrm{ab}} \dotimes{\Z} \Bbbk \arrow[r, "\sim"] & \Hm_1(G, \Bbbk) \\
		(g \bmod G_{\mathrm{der}}) \otimes 1 \arrow[mapsto, r] & {[g] \;\text{modulo the relations}},
	\end{tikzcd}\]
	where $G_{\mathrm{ab}}:=G/G_{\mathrm{der}}$ is the abelianization of $G$
	and $G_{\mathrm{der}}$ is its derived subgroup, as in
	\cite[Definition 4.7.2]{Li1}. Readers who know the topological background
	of group homology will recognize the connection with the Hurewicz theorem
	in homology theory.

	It is natural also to ask for the composite of the isomorphism
	$\Hm_1(G,\Bbbk)\rightiso\mathfrak{I}/\mathfrak{I}^2$ from
	Example~\ref{eg:group-H1} with the isomorphism above. We claim that this
	composite is given explicitly by
% segment-id: o014.li2.chapter6.group-homology-cohomology.g005
	\[\begin{tikzcd}[row sep=tiny]
		G_{\mathrm{ab}} \otimes \Bbbk \arrow[r, "\sim"] & \mathfrak{I}/\mathfrak{I}^2 \\
		(g \bmod G_{\mathrm{der}}) \otimes 1 \arrow[mapsto, r] & 1_G - g \;\bmod \mathfrak{I}^2.
	\end{tikzcd}\]

	Why is this so? The isomorphism in Example~\ref{eg:group-H1} comes from
	the map $\Hm_1(G,\Bbbk)\to\Hm_0(G,\mathfrak{I})\simeq
	\mathfrak{I}/\mathfrak{I}^2$ induced by the short exact sequence
	$0\to\mathfrak{I}\to\Bbbk[G]\to\Bbbk\to0$. Write down the following
	commutative diagram with exact rows:
% segment-id: o014.li2.chapter6.group-homology-cohomology.g006
	\[\begin{tikzcd}
		& & {(g)^\dagger \otimes 1_G} \arrow[phantom, d, "\in" description, sloped] \arrow[mapsto, r] & {[g] = (g)^\dagger \otimes 1} \arrow[phantom, d, "\in" description, sloped] & \\
		0 \arrow[r] & C_1(G, \mathfrak{I}) \arrow[r] \arrow[d] & C_1(G, \Bbbk[G]) \arrow[r] \arrow[d, "{\partial'_1 \otimes \identity_{\Bbbk[G]}}"] & C_1(G, \Bbbk) \arrow[r] \arrow[d, "{\partial'_0 \otimes \identity_{\Bbbk[G]} = 0}"] & 0 \\
		0 \arrow[r] & C_0(G, \mathfrak{I}) \arrow[r] & C_0(G, \Bbbk[G]) \arrow[r] & C_0(G, \Bbbk) \arrow[r] & 0 ,
	\end{tikzcd}\]
	The map $\partial'_1\otimes\identity_{\Bbbk[G]}$ sends
	$(g)^\dagger\otimes1_{\Bbbk[G]}$ to
	$1_G-g\in C_0(G,\mathfrak{I})=\mathfrak{I}$, thereby determining an
	element of $\mathfrak{I}/\mathfrak{I}^2\simeq
	\Hm_0(G,\mathfrak{I})$.\footnote{Translator's correction: the source uses
	$C^0$ and $\Hm^0$, but the diagram and the connecting morphism being
	computed are homological, so the correct indices are $C_0$ and $\Hm_0$.}
	This is precisely the construction in the snake lemma.
\end{example}

% segment-id: o014.li2.chapter6.group-homology-cohomology.r001
\begin{remark}\label{rem:Gmod-k}
	The form of the standard cochain complex (respectively standard chain
	complex) shows that, as an additive group, $\Hm^n(G,M)$ (respectively
	$\Hm_n(G,M)$) is completely determined by the additive structure of $M$;
	multiplication by elements of $\Bbbk$ only supplies its $\Bbbk$-module
	structure. Thus the case $\Bbbk=\Z$ already captures the essence of group
	cohomology (respectively homology), although allowing general coefficients
	causes no additional difficulty.
\end{remark}

% segment-id: o014.li2.chapter6.group-homology-cohomology.r002
\begin{remark}[Normalized version]\label{rem:nor-cochain}
	\index{standard cochain complex!normalized}
	\index[sym1]{CGMbar@$\overline{C}(G, M)$}
	To compute $\Hm^n(G,M)$ and $\Hm_n(G,M)$, one may instead use the
	normalized free resolution $\overline{\mathsf{L}}\to\Bbbk$ supplied by
	Proposition~\ref{prop:trivial-mod-nor-resolution}. Denote the corresponding
	cochain and chain complexes by $(\overline{C}^n(G,M))_n$ and
	$(\overline{C}_n(G,M))_n$, respectively, and call them the normalized
	standard cochain complex and normalized standard chain complex. They are,
	respectively, a subcomplex and a quotient complex of the original
	versions. For example, in the cohomological case, the description of
	$\mathsf{L}_n$ in \eqref{eqn:nor-cochain} gives
% segment-id: o014.li2.chapter6.group-homology-cohomology.q024
	\[ \overline{C}^n(G, M) = \left\{\begin{array}{r|l}
		f \in C^n(G, M) & \exists 1 \leq k \leq n, \; g_k = 1_G \\
		& \implies f(g_1, \ldots, g_n) = 0
	\end{array}\right\}. \]
\end{remark}
